\section*{Chapter \ref{ch:b2:insurance-and-risk-sharing} --- Insurance and Risk Sharing}

\begin{solution}{exo:b2:insurance-and-risk-sharing:1}
The United States lies above the line ($1.97 > 1.76$): its consumption moved more than its output.
% figdata: volatility (USA, JPN)
Japan lies below ($1.50 < 2.05$). Under full sharing with identical CRRA utilities, every country's
consumption grows at the world's rate, so every country has the world's consumption volatility, 1.79
% figdata: summary (world_sd_cons 1.79)
points: the points would lie on a horizontal line at 1.79, below the dotted line for every country
whose GDP is more volatile than that (all but 8).
\end{solution}

\begin{solution}{exo:b2:insurance-and-risk-sharing:2}
Fair premium $0.005 \times \num{40000} = \dal{200}$. \omterm{def:b2:insurance-and-risk-sharing:loading}{Loading} $260/200 - 1 = 30$ per cent.
\end{solution}

\begin{solution}{exo:b2:insurance-and-risk-sharing:3}
% code: exercises/solutions.py ex_contingent
$z_G = \num{50000} - 0.12\chi$, $z_B = \num{30000} + 0.88\chi$. Budget line slope $-0.88/0.12 = -7.33$;
fair-odds slope $-0.9/0.1 = -9$; \omterm{def:b2:insurance-and-risk-sharing:loading}{loading} $0.12/0.1 - 1 = 20$ per cent. On the certainty line,
$\num{50000} - 0.12\chi = \num{30000} + 0.88\chi$, so $\chi = \num{20000}$ (full cover) and $z_G = z_B
= \dal{47600}$.
\end{solution}

\begin{solution}{exo:b2:insurance-and-risk-sharing:4}
Rejections: a household's consumption falls in a year when its own income falls (an illness, a
lay-off, a hailstorm on its fields only) while the village's average consumption does not; or the gaps
between households' consumptions widen after shocks to their own incomes and do not close. Not a
rejection: every household's consumption falls in a drought year. A drought is \omterm{def:b2:insurance-and-risk-sharing:aggregate-risk}{aggregate risk}, which
no sharing among the villagers can remove.
\end{solution}

\begin{solution}{exo:b2:insurance-and-risk-sharing:5}
% code: exercises/solutions.py ex_cover
With the formula of \cref{ex:b2:insurance-and-risk-sharing:loading}: $\chi^{*} = (0.1 \times 0.88
\times \num{50000} - 0.9 \times 0.12 \times \num{30000})/(0.12 \times 0.88) = (\num{4400} -
\num{3240})/\num{0.1056} = \dal{10985}$, about half of full cover. Wealth $\num{50000} - 0.12 \times
\num{10985} = \dal{48682}$ with no loss and $\num{30000} + 0.88 \times \num{10985} = \dal{39667}$ after
a loss.
\end{solution}

\begin{solution}{exo:b2:insurance-and-risk-sharing:6}
A higher price means a higher \omterm{def:b2:insurance-and-risk-sharing:loading}{loading}, so cover falls. At $\psi = 0.15$: $\chi^{*} = (0.1 \times 0.85
\times \num{50000} - 0.9 \times 0.15 \times \num{30000})/(0.15 \times 0.85) = 200/0.1275 =
\dal{1569}$. It buys nothing once $0.1(1 - \psi)\num{50000} \le 0.9\psi \times \num{30000}$, that is
$\psi \ge \num{5000}/\num{32000} = 0.15625$, a \omterm{def:b2:insurance-and-risk-sharing:loading}{loading} of $56.25$ per cent. From a 20 to a 50 per cent
\omterm{def:b2:insurance-and-risk-sharing:loading}{loading}, cover falls by six sevenths.
\end{solution}

\begin{solution}{exo:b2:insurance-and-risk-sharing:7}
% code: exercises/solutions.py ex_borch
By \cref{ex:b2:insurance-and-risk-sharing:cara}, $z_1 = \tfrac{a_2}{a_1 + a_2}\Omega + k =
\tfrac13\Omega + k$. The expected harvest is \num{80000}, so $k = \num{30000} - \num{80000}/3 =
\dal{3333}$. Farmer 1 gets \dal{36667} or \dal{23333}; farmer 2 gets \dal{63333} or \dal{36667}. Of the
\dal{40000} swing in the harvest, farmer 2 bears two thirds: he is half as risk-averse, so he has twice
the risk tolerance.
\end{solution}

\begin{solution}{exo:b2:insurance-and-risk-sharing:8}
% code: exercises/solutions.py ex_six_countries
Standard deviations (GDP, consumption) and ratio: United States $1.76$, $1.97$, $1.12$; Japan $2.05$,
$1.50$, $0.73$; Korea $3.23$, $4.40$, $1.36$; France $2.22$, $1.85$, $0.83$; Brazil $2.65$, $3.13$,
$1.18$; India $2.89$, $2.83$, $0.98$. The prediction is a consumption volatility equal to the world's,
1.79 points, and a ratio below one: France, Japan and the United States come closest in level, France
and Japan also in the ratio. Korea is furthest away.
\end{solution}

\begin{solution}{exo:b2:insurance-and-risk-sharing:9}
Full sharing (identical CRRA utilities) puts every point at a correlation of $1$ for consumption:
along the top edge of the chart. Measured correlations can fall short without any failure of
sharing because the series is private consumption only (public services are left out), because much of
consumption is non-traded (housing, local services) and its relative price moves with local
conditions, because consumption is measured with error, and because tastes differ from year to year.
To separate these from a lack of sharing, look at the gap where measurement is good and for traded
goods only; or ask whether the gap is larger in countries that hold less foreign wealth. Measurement
error does not explain why holdings of foreign assets should matter.
\end{solution}

\begin{solution}{exo:b2:insurance-and-risk-sharing:10}
Borch: $c_1^{-\gamma_1}/c_2^{-\gamma_2} = \kappa$, so $\gamma_1\ln c_1 - \gamma_2\ln c_2 = -\ln\kappa$
in every state. Across states, $\gamma_1\,\dd\ln c_1 = \gamma_2\,\dd\ln c_2$; both consumptions rise
with the aggregate (\cref{cor:b2:insurance-and-risk-sharing:aggregate-only}), and household 1's grows
$\gamma_2/\gamma_1 > 1$ times as fast in percentage terms, so its share rises with the aggregate. In the
regression, household 1's consumption moves more than one for one with average consumption ($\beta_1
> 1$), household 2's less, and own income still adds nothing ($\beta_2 = 0$). So $\beta_1 \neq 1$ does
not reject full sharing when attitudes to risk differ; the test that survives is $\beta_2 = 0$.
\end{solution}

\begin{solution}{exo:b2:insurance-and-risk-sharing:11}
% code: exercises/solutions.py ex_pooling
Each household bears $\frac1n\sum_j\tilde x_j$, and its cost $\Gamma$ solves $u(z - \Gamma) =
\mathbb{E}[u(z - \frac1n\sum_j \tilde x_j)]$, which gives $e^{a\Gamma} = \mathbb{E}[e^{a\sum_j \tilde x_j/n}] = \bigl(1 - \mu +
\mu e^{a\Lambda/n}\bigr)^n$ by independence. Numbers: \dal{1585.7} for $n = 1$, \dal{1046.2} for $n =
10$, \dal{1004.5} for $n = 100$, against an expected loss of \dal{1000}. Write the cost as $K(s)/s$ with
$s = a/n$ and $K(s) = \ln\mathbb{E}[e^{s\tilde x}]$. $K$ is convex with $K(0) = 0$, so $K(s)/s$ rises
with $s$: the cost falls as $n$ rises, towards $K'(0) = \mathbb{E}[\tilde x] = \mu\Lambda$. With
perfectly correlated losses, everyone loses together, each household bears its full loss in that state,
and the cost stays \dal{1585.7} whatever $n$.
\end{solution}

\begin{solution}{exo:b2:insurance-and-risk-sharing:12}
% code: exercises/solutions.py ex_portfolio
$\theta^{*} = 0.05/(3 \times 0.18^2) = 0.514$. Under CRRA, the share is the same at every wealth, so
the amount doubles when wealth doubles. Under CARA with normal returns, the amount $\bar\xi/(a\sigma^2)$
is fixed, so the share halves. Data in which the rich and the poor hold the same share reject CARA;
data in which they hold the same amount reject CRRA; a share that rises with wealth rejects both, and
calls for relative \omterm{def:b2:expected-utility:risk-aversion}{risk aversion} that falls with wealth.
\end{solution}

\begin{solution}{pb:b2:insurance-and-risk-sharing:1}
% code: exercises/solutions.py problem, problem_bb_cover
\textbf{1.} A boat is good with probability $0.4 + 0.2 = 0.6$; both are good with probability $0.4 >
0.6^2 = 0.36$. The catches are not independent: they are positively correlated (a common sea).

\textbf{2.} Mean $0.6 \times \num{30000} + 0.4 \times \num{10000} = \kst{22000}$; standard deviation
$\sqrt{0.6 \times 0.4} \times \num{20000} = \kst{9798}$.

\textbf{3.} Covariance $\num{20000}^2(0.4 - 0.36)$, variance $\num{20000}^2 \times 0.24$: correlation
$0.04/0.24 = 0.167$.

\textbf{4.} \kst{60000} with probability $0.4$, \kst{40000} with probability $0.4$ (one boat good),
\kst{20000} with probability $0.2$.

\textbf{5.} $0.6\ln\num{30000} + 0.4\ln\num{10000} = 9.8695$; \omterm{def:b2:expected-utility:certainty-equivalent}{certainty equivalent} \kst{19332}.

\textbf{6.} The remainder $z_1 - \Omega/2$ is $0$ when both seasons are alike, $+\num{10000}$ when only
boat 1 is good and $-\num{10000}$ when only boat 2 is; household 2's is the opposite, so the two add up
to zero. Standard deviations: \kst{7483} for $\Omega/2$, \kst{6325} for the remainder (variances
$5.6$ and $4.0$ times $10^7$, which add up to $9.6 \times 10^7$, the variance of the catch).

\textbf{7.} Both have CRRA utility with $\gamma = 1$; with $n = 2$,
\cref{prop:b2:insurance-and-risk-sharing:constant-shares} gives $z_1(s) = k\,\Omega(s)$ with a constant
$k$, and $z_2 = (1 - k)\Omega$.

\textbf{8.} \kst{30000}, \kst{20000}, \kst{20000}, \kst{10000} in the four states: mean \kst{22000},
unchanged; expected utility $0.4\ln\num{30000} + 0.4\ln\num{20000} + 0.2\ln\num{10000} = 9.9270$;
\omterm{def:b2:expected-utility:certainty-equivalent}{certainty equivalent} \kst{20477}.

\textbf{9.} $\num{20477} - \num{19332} = \kst{1145}$ a year for each household.

\textbf{10.} The variance falls from $9.6$ to $5.6$ times $10^7$, by 41.7 per cent. The idiosyncratic
part, the remainder of question 6, is removed entirely: 100 per cent.

\textbf{11.} \Cref{ex:b2:insurance-and-risk-sharing:cara} with $a_1 = 2a_2$: skipper 1 bears
$a_2/(a_1 + a_2) = \tfrac13$ of the aggregate, $z_1 = \Omega/3 + k$, and $k = \num{22000} -
\num{44000}/3 = \kst{7333}$. Skipper 1 gets \kst{27333} when both are good, \kst{20667} when one is,
\kst{14000} when both are bad; skipper 2 gets the rest: \kst{32667}, \kst{19333}, \kst{6000}.

\textbf{12.} What remains is \omterm{def:b2:insurance-and-risk-sharing:aggregate-risk}{aggregate risk}: in a bad year for both boats, there is less to share.
Only someone whose income does not fall in the same seasons can take it: an insurer that pools
fisheries of many coasts, or investors elsewhere, at a price.

\textbf{13.} A household's bad season has probability $0.4$: the fair price is $0.4$ per kest of cover,
the loaded price $0.48$.

\textbf{14.} $z_0 = \num{30000}$, $\Lambda = \num{20000}$, $\mu = 0.4$, $\psi = 0.48$: $\chi^{*} =
(0.4 \times 0.52 \times \num{30000} - 0.6 \times 0.48 \times \num{10000})/(0.48 \times 0.52) =
\num{3360}/\num{0.2496} = \kst{13462}$.

\textbf{15.} $\num{30000} - 0.48 \times \num{13462} = \kst{23538}$ in a good season, $\num{10000} +
0.52 \times \num{13462} = \kst{17000}$ in a bad one; \omterm{def:b2:expected-utility:certainty-equivalent}{certainty equivalent} \kst{20666}.

\textbf{16.} The insurer, by \kst{189} a year. The cooperative removes only the \omterm{def:b2:insurance-and-risk-sharing:idiosyncratic-risk}{idiosyncratic risk}; the
insurer, which pools many coasts, takes some of the \omterm{def:b2:insurance-and-risk-sharing:aggregate-risk}{aggregate risk} too, and here that is worth more than
the \omterm{def:b2:insurance-and-risk-sharing:loading}{loading}. The insurer expects to earn $(0.48 - 0.4) \times \num{13462} = \kst{1077}$ a year on the
contract.

\textbf{17.} The price of cover paid in the state where both are bad is $1.2 \times 0.2 = 0.24$.
Expected utility $0.4\ln(\num{30000} - 0.24\chi) + 0.4\ln(\num{20000} - 0.24\chi) + 0.2\ln(\num{10000}
+ 0.76\chi)$ is largest at $\chi^{*} = \kst{9391}$; \omterm{def:b2:expected-utility:certainty-equivalent}{certainty equivalent} \kst{21072}.

\textbf{18.} Cooperative with cover of the bad state for both (\kst{21072}), insurer alone (\kst{20666}),
cooperative alone (\kst{20477}), nothing (\kst{19332}). A cover paid on a household's own bad season
weakens the reward for the skipper's effort (\cref{ch:b2:moral-hazard}); sharing in the cooperative
does too, since each household keeps only half of its own catch, but the members can watch each
other's work more closely than an insurer can. A cover paid only when both boats fail depends little
on any one skipper's effort.

\textbf{19.} $\beta_1 = 1$ and $\beta_2 = 0$: with identical CRRA utilities every household's
consumption is a constant share of the total
(\cref{prop:b2:insurance-and-risk-sharing:constant-shares}), so its growth is the group's, whatever its
own income did.

\textbf{20.} $\hat\beta_1 = 0.853$ (standard error $0.034$), $\hat\beta_2 = 0.143$ ($0.006$).
% figdata: coop_regression

\textbf{21.} $t = 22.6$ (the table divides unrounded figures; $0.143/0.006$ gives about 24), far
above $1.96$: $\beta_2 = 0$ is rejected.

\textbf{22.} $t = (0.853 - 1)/0.034 = -4.3$: $\beta_1 = 1$ is rejected too. The sum is $0.996$. If a
household keeps the share $\kappa$ of its own shock, $\Delta\ln c_{it} \approx \Delta m^{c}_t +
\kappa\,\Delta m^{o}_{it}$, where $m^{c}_t$ is the common part and $m^{o}_{it}$ the household's own
part of log income. Since $\Delta\ln m_{it} = \Delta m^{c}_t + \Delta m^{o}_{it}$ and
$\Delta\ln\bar c_t \approx \Delta m^{c}_t$, this is $(1 - \kappa)\Delta\ln\bar c_t + \kappa\,\Delta\ln m_{it}$: the two coefficients add up to
one.

\textbf{23.} Households keep about 14 per cent of their own income shocks; the cooperative insures about
86 per cent of idiosyncratic income risk.

\textbf{24.} Error in measured income biases $\hat\beta_2$ towards zero, so the cooperative would look
as if it insured more than 86 per cent. Error in consumption only widens the standard errors.

\textbf{25.} The two-boat cooperative is worth \kst{1145} a year to each household, by removing all of
its \omterm{def:b2:insurance-and-risk-sharing:idiosyncratic-risk}{idiosyncratic risk}; the larger cooperative insures about 86 per cent of its members' idiosyncratic
income risk ($1 - \hat\beta_2$).
\end{solution}
